% VI-B Runtime Configuration Calibration
\subsection{Runtime Configuration Calibration}\label{sec:runtime-configuration-calibration}
\paragraph{Local concurrency.}
We calibrated the Local in-flight limit with the full Local inference pipeline under the eight-stream, 240-FPS offered workload. Calibration05 was the preregistered selection dataset. The two-repeat mean completion rates for $C_L\in\{3,4,5,6\}$ were 232.800, 234.858, 235.475, and 234.767 FPS, respectively. None sustained the full offered load. The preregistered fallback therefore selected the smallest tested concurrency within 2\% of the maximum observed throughput. The maximum was 235.475 FPS, giving a threshold of 230.766 FPS; $C_L=3$ achieved 232.800 FPS and was formally selected. Additional in-flight concurrency yielded little aggregate-throughput gain while increasing per-frame service time in these measurements. The subsequent $K=1\ldots8$ Local-only sweep characterizes the selected configuration; it did not select $C_L$ retrospectively.

\paragraph{Edge concurrency.}
The historical Edge-concurrency rule compared repeat-matched improvements in overall timely-inference ratio (TIR) for temporally concentrated placement, using a preregistered absolute threshold of 0.10. The two observed improvements were 0.138 and 0.025, so the formal verdict was \texttt{EDGE\_C\_REPEAT\_AMBIGUOUS}; that rule selected no $C_E$. After this result and before Grid03 measurement, we prospectively fixed $C_E=2$ as a conservative runtime configuration. In the concentrated stress condition, increasing $C_E$ from one to two reduced observed server queue-wait p95 from approximately 10.8--18.5\,ms to 0.8--1.2\,ms. This largely, but not completely, removed the observed single-worker server queueing and made $C_E=2$ a more conservative runtime for evaluating temporal placement. We do not claim that $C_E=2$ maximizes Edge throughput.

$C_E=2$ is conservative with respect to our temporal-placement hypothesis because increasing Edge concurrency improved the temporally concentrated baseline and reduced rather than amplified the placement gap in the robustness measurements. Even under $C_E=2$, temporally dispersed placement retained TIR advantages of +0.161 and +0.241 over concentrated placement in the two repeats. These gaps are supporting robustness evidence, not the criterion used to choose the forward Edge runtime. The final primary runtime is $C_L=3$, $C_E=2$, with TensorRT batch size $B=1$. The implementation labels \texttt{ALIGNED} and \texttt{STAGGERED} denote temporally concentrated and temporally dispersed assignment placement, respectively; camera-source synchronization was not manipulated.
